\section{Research Roadmap and Limits}
\label{sec:open-problems}

\subsection{Physically Grounded Attacks}

Physical grounding is essential because VLA attacks ultimately claim effects on robot motion, tool use, environment state, or human proximity. Current work has moved beyond digital-only examples: several patch, texture, backdoor, and embodied red-team studies include simulation or real-robot demonstrations. Progress now depends less on simply adding physical trials than on making those trials more informative: calibrated cameras, lighting and viewpoint variation, object mass and friction, controller frequency, safety filters, time-to-violation, and recovery opportunities should be part of the evaluated claim.

\subsection{HIL and Real-Robot Recovery}

Recovery is the most underdeveloped robotics question. Simulation gives scale, and full physical deployment gives realism, but hardware-in-the-loop can expose real sensors, middleware timing, safety filters, controller latency, and emergency-stop logic without requiring unsafe full deployment. Future studies should show when an attack first changes representation, whether the controller clips or rejects it, whether the robot recovers after new observations, and whether human correction changes the outcome. A VLA attack that fails after one replan and a VLA attack that survives repeated recovery attempts are different risks.

\subsection{Cross-Embodiment and Action-Representation Transfer}

Transfer matters because VLA systems are defined by embodiment as much as by model name. Existing work has begun to test transfer across prompts, models, viewpoints, and simulation-to-real settings, but transfer across action representations is still thin. A prompt attack may transfer across language-conditioned token policies but fail against a flow policy; a patch may transfer across visual encoders but not camera placements; a chunk-drift backdoor may be irrelevant to a one-step replanner. Future suites should vary tokenizer, chunk length, temporal ensemble, diffusion/flow sampler, waypoint controller, robot morphology, camera geometry, and low-level safety envelope.

\subsection{Memory, Tool, Middleware, and HRI Attacks}

Deployed VLA robots will not be stateless policies. They will retrieve maps, object memories, user preferences, logs, tool outputs, and social cues. Adjacent evidence from RAG, software agents, ROS/ROS2, and HRI shows credible routes for corruption \cite{zou2025poisonedrag,debenedetti2024agentdojo,deng2022insecurityros2,booth2017piggybacking}, but direct VLA evidence remains thin. Future work should measure robot-level effects: unauthorized skill use, corrupted maps, persistent user preference changes, unsafe tool arguments, privacy leakage, and recovery after correction.

\subsection{Privacy as a First-Class VLA Attack Objective}

Robot actions can reveal training data, user routines, scenes, and task histories. Membership inference work shows that VLA actions and trajectories can leak training membership \cite{peng2026miavla,luan2026vlaleaks}. The field now needs privacy tests beyond membership: property inference, dataset provenance, user routine leakage, private-scene memorization, and cross-session log leakage. These are direct VLA security questions even when no robot collision or task failure occurs.

\subsection{Adaptive Attacks and Responsible Artifacts}

First-generation defenses will use prompt hierarchy, source labels, retrieval filters, tool permissions, runtime monitors, action shields, control barrier functions, and human approvals. Attack research must therefore include adaptive settings: what the attacker knows, what feedback is available, whether queries are budgeted, and whether benign utility is preserved. At the same time, VLA attack artifacts can be unusually transferable to physical systems. Benchmarks should separate public scoring code and sanitized examples from high-risk payloads, physical exploit details, or deployment-specific trigger recipes.

\subsection{Limitations of This Review}

This review fixes its evidence set at the 16 June 2026 search date, so later arXiv releases, renamed project pages, and venue updates are outside scope. Google Scholar and Semantic Scholar ranking can bias discovery toward highly visible papers, while targeted-name searches can miss attacks that use different terminology. We searched English-language sources and may miss non-English reports. Several direct VLA attacks are recent preprints or workshop papers; their claims may change after peer review, and project-page or repository information is weaker evidence than archival text.

The supplementary materials document inclusion choices and supporting notes, but they do not reproduce experiments, inspect every codebase, or verify every ASR number. Some judgments remain partly subjective, especially for target representations, transfer, and validation details. The descriptive observations should therefore be read as a map of an emerging literature, not as a meta-analysis; we do not pool ASR values across attack objectives, denominators, validation tiers, or robot platforms.
